\section{Experimental Design}

The full study should evaluate long-context quality, calibration, memory use, and throughput on standard benchmarks. This draft reports only a small synthetic retrieval pilot, described in Section~\ref{sec:pilot-results}. The broader protocol below is retained to make the intended study reproducible and to prevent accidental cherry-picking.

\subsection{Research questions}

\begin{enumerate}[leftmargin=*]
    \item \textbf{RQ1: Quality-efficiency tradeoff.} At the same memory budget, does \uckv{} preserve task quality better than non-calibrated cache compression baselines?
    \item \textbf{RQ2: Calibration.} Are the uncertainty and compression-risk estimates calibrated across datasets, context lengths, and answer types?
    \item \textbf{RQ3: Selective risk control.} Does lowering the risk tolerance $\epsilon$ monotonically reduce observed compression degradation while increasing memory use?
    \item \textbf{RQ4: Failure modes.} When \uckv{} fails, are errors associated with long-range retrieval, middle-position evidence, multi-hop reasoning, high entropy, or poor calibration?
\end{enumerate}

\subsection{Models}

Initial open-weight model candidates are Llama-2-7B/13B \citep{touvron2023llama}, Mistral-7B \citep{jiang2023mistral}, and Llama-3-8B if licensing and local hardware permit \citep{grattafiori2024llama}. The first implementation should start with one 7B-class model to stabilize instrumentation before scaling to additional model families.

\subsection{Datasets}

Long-context evaluation should use LongBench \citep{bai2024longbench}, RULER \citep{hsieh2024ruler}, and lost-in-the-middle style key-value retrieval placements \citep{liu2024lost}. Reliability-oriented evaluation should include TruthfulQA \citep{lin2022truthfulqa}, TriviaQA \citep{joshi2017triviaqa}, Natural Questions \citep{kwiatkowski2019natural}, and selected MMLU subsets \citep{hendrycks2021mmlu}. Language-modeling sanity checks can use WikiText-2 and C4-style held-out text if licensing and preprocessing are handled consistently.

\subsection{Baselines}

Baselines should include:

\begin{itemize}[leftmargin=*]
    \item full KV cache;
    \item sliding window and sink-plus-window retention;
    \item H2O-style heavy-hitter retention \citep{zhang2023h2o};
    \item Scissorhands-style persistent importance \citep{liu2023scissorhands};
    \item SnapKV-style prompt compression \citep{li2024snapkv};
    \item PyramidKV or Ada-KV style adaptive allocation where implementation is feasible \citep{cai2024pyramidkv,feng2024adakv};
    \item KIVI or KVQuant as optional quantization baselines or compositional extensions \citep{liu2024kivi,hooper2024kvquant}.
\end{itemize}

\subsection{Metrics}

Task quality metrics should follow benchmark conventions: exact match, F1, ROUGE, multiple-choice accuracy, retrieval accuracy, and perplexity where appropriate. Reliability metrics should include ECE, adaptive ECE, Brier score, negative log-likelihood, selective risk-coverage curves, and area under the risk-coverage curve. Efficiency metrics should include peak GPU memory, average KV cache bytes per request, tokens per second, time to first token, decode latency, and maximum batch size under a fixed memory limit.

\subsection{Pilot setup used in this draft}

The pilot uses controlled key-value retrieval prompts generated by the synthetic benchmark script in \texttt{experiments/}. Each prompt contains distractor text and one key-value pair, and the model is asked to return the value for a named key. Prompts vary across four context lengths (64, 128, 192, and 256 approximate words) and three answer positions (early, middle, and late). We use answer containment in the generated continuation as a simple task-level metric.

The final reported pilot uses SmolLM2-135M-Instruct with the tokenizer chat template, greedy decoding, and a 24-token generation budget. A main sweep uses 144 prompts from seed 29 and sink-plus-window candidates with four protected sink tokens and recent windows $\{32,64,128,\allowbreak 192,256,320,384\}$. The risk model is a logistic classifier trained on replay features to predict whether a candidate compression action causes step-level KL divergence above 0.01 relative to full-cache decoding. We tune \uckv{} risk tolerance on this main sweep and then run two confirmation seeds, 31 and 37, with 96 prompts each.

For each full replay run, the script logs full-cache greedy continuations, fixed-window replays, step-level KL divergence from the full-cache next-token distribution, top-1 mismatch, retained-token counts, logistic compression-risk calibration metrics, \uckv{} selections, and free-running compressed generations. Tolerance sweeps can reuse a fixed-window replay run and rerun only \uckv{} free-generation; those runs report real task generations but mark replay metrics as fixed-candidate simulations. The confirmation seeds report actual replay metrics for all policies.

\subsection{Reporting templates for full study}

The full-benchmark templates below remain placeholders. They define the intended reporting shape for future LongBench/RULER and reliability runs.

\begin{table}[h]
\centering
\caption{Quality-efficiency comparison template. Values are intentionally omitted in this draft.}
\begin{tabular}{lcccc}
\toprule
Method & KV budget & Task score & Peak memory & Decode tok/s \\
\midrule
Full KV & 100\% & TBD & TBD & TBD \\
Baseline compression & TBD & TBD & TBD & TBD \\
\uckv{}-Budget & TBD & TBD & TBD & TBD \\
\uckv{}-Score & TBD & TBD & TBD & TBD \\
\uckv{}-Full & TBD & TBD & TBD & TBD \\
\bottomrule
\end{tabular}
\end{table}

\begin{table}[h]
\centering
\caption{Calibration and selective-risk template. Values are intentionally omitted in this draft.}
\begin{tabular}{lcccc}
\toprule
Method & ECE & Brier & Risk at coverage & AU-RC \\
\midrule
Token entropy & TBD & TBD & TBD & TBD \\
Calibrated uncertainty & TBD & TBD & TBD & TBD \\
Compression-risk model & TBD & TBD & TBD & TBD \\
\bottomrule
\end{tabular}
\end{table}

\subsection{Ablations}

The ablation plan should vary uncertainty features, calibration method, risk tolerance $\epsilon$, sink size, recent window size, score decay $\gamma$, uncertainty weight $\eta$, per-head versus shared retained sets, and prompt-only versus online compression. For each ablation, report both quality and reliability metrics, not only memory savings.
